\section{\S 2.1 Explain the Molecular Structure by VSEPR and Hybrid Orbital Theory}\label{sect:VSEPR和杂化轨道理论}

    \subsection{价层电子对互斥理论（VSEPR）}\label{subsec:VSEPR}

        \begin{itemize}
            \item \makebox[10em][l]{\textbf{\uline{V}alence}}价
            \item \makebox[10em][l]{\textbf{\uline{S}hell}}层
            \item \makebox[10em][l]{\textbf{\uline{E}lectron}}电子
            \item \makebox[10em][l]{\textbf{\uline{P}air}}对
            \item \makebox[10em][l]{\textbf{\uline{R}epulsion}}排斥
        \end{itemize}

        能够解释、判断简单共价分子和（含共价键的）离子的几何构型。

        主要适用于 \(AB_n\) 型分子或离子

        \begin{itemize}
            \item A——中心原子
            \item B——配位原子
            \item n——配位原子个数
        \end{itemize}

        \subsubsection{理论基本观点}\label{subsubsec:VSEPR基本观点}

            中心原子的价层电子对由于相互间的排斥作用，尽可能彼此远离，使电子对之间的静电斥力最小，体系能量最低，分子最稳定。

            价层电子对包括：成键 \(\upsigma\) 电子对和孤电子对

        \subsubsection{价层电子对的排列方式}\label{subsubsec:价层电子对排列方式}

            \begin{center}
            \begin{tabular}{|c|c|c|c|c|c|c|}
                \hline
                \noindent\textbf{价层电子对数} & \textbf{2对} & \textbf{3对} & \textbf{4对} & \textbf{5对} & \textbf{6对} \\
                \hline
                \makecell{VSEPR \\ 理论模型} & 直线形 & 平面正三角形 & 正四面体形 & & \\
                \hline
                \makecell{中心原子 \\ 杂化方式} & sp & sp$^2$ & sp$^3$ & & \\
                \hline
            \end{tabular}
            \end{center}

        \subsubsection{确定价层电子对数}\label{subsubsec:确定价层电子对数}

            价层电子对数 =
            \[\underset{\text{成键}\upsigma\text{电子对数}}{n} + (\underset{\text{孤电子对数}}{\text{A的价电子数}} - \overset{\makebox[0pt]{\footnotesize 代数值}}{\text{离子电荷}} - n \times \text{B的单电子数})\div 2\]

            或：\quad   \(\begin{cases} - \text{ 正离子电荷数} \\ + \text{ 负离子电荷数} \end{cases}\)

            \begin{xmp}{价层电子对数计算}
                计算下列分子或离子的价层电子对数。
                
                \cemh{CO2}：\(2 + 0 = 2\)对\qquad \cemh{PCl3}：\(3 + 1 = 4\)对\qquad \cemh{SO2}：\(2 + 1 = 3\)对
                \qquad \cemh{NH4+}：\(4 + 0 = 4\)对\qquad \cemh{CO3^{2-}}：\(3 + 0 = 3\)对
            \end{xmp}

        \subsubsection{实例分析}\label{subsubsec:VSEPR实例分析}

            \begin{center}
            \small
            \begin{tabular}{|c|c|c|c|c|c|}
                \hline
                \noindent\textbf{\makecell{价层电子\\对数}} & \textbf{\makecell{VSEPR\\模型}} & \textbf{\makecell{成键\(\upsigma\)\\电子对数}} & \textbf{\makecell{孤电子\\对数}} & \textbf{实例} & \textbf{分子几何构型} \\
                \hline
                \multirow{1}{*}{2} & 直线形 & 2 & 0 & \makecell{\cemh{BeCl2}, \cemh{CO2}} & 直线形 \\
                \hline
                \multirow{2}{*}{3} & 正三角形 & 3 & 0 & \makecell{\cemh{BF3}, \cemh{CO3^{2-}}} & 正三角形 \\
                \cline{3-6}
                & 三角形 & 2 & 1 & \cemh{SO2} & \makecell{角形 \\\footnotesize（折线形、V字形）} \\
                \hline
                \multirow{3}{*}{4} & 正四面体形 & 4 & 0 & \makecell{\cemh{CCl4}, \cemh{CH4},\\\cemh{NH4+}} & 正四面体形 \\
                \cline{3-6}
                & 四面体形 & 3 & 1 & \makecell{\cemh{NH3}, \cemh{PCl3},\\\cemh{PH3}} & 三角锥形 \\
                \cline{3-6}
                & 四面体形 & 2 & 2 & \makecell{\cemh{H2O}, \cemh{H2S}} & \makecell{角形 \\\footnotesize（折线形、V字形）} \\
                \hline
            \end{tabular}
            \end{center}

            \noindent\textbf{总结}
            \begin{itemize}
                \item \textcolor{blue}{先计算出价层电子对总数，确定VSEPR模型；}
                \item \textcolor{blue}{若孤电子对数 = 0，则分子的空间构型即为VSEPR模型；}
                \item \textcolor{blue}{若孤电子对数 \(\neq\) 0，则从VSEPR模型扣除相应的顶点数，判断剩余原子组成的空间构型。}
            \end{itemize}

        \subsubsection{影响键角的因素}\label{subsubsec:影响键角的因素}

        \vspace{-\baselineskip}
            \begin{align*}
                \cemh{CH4} \text{键角：} & \ang{109;28} \\
                \cemh{NH3} \text{键角：} & \ang{107;18} \\
                \cemh{H2O} \text{键角：} & \ang{104.5}
            \end{align*}

            都是4对价层电子对，相同的VSEPR模型，为什么键角不同？

            \textcolor{red}{斥力大小}：lp-lp \(>\) lp-bp \(>\) bp-bp\quad \(\longleftarrow\)\textcolor{red}{为什么？}

            （lp：孤电子对 lonely pair；bp：成键电子对 bonding pair）

            \textcolor{blue}{\(\therefore\) 键角：\cemh{H2O} \(<\) \cemh{NH3} \(<\) \cemh{CH4}}

            \begin{xmp}{分子构型分析}
                分析下列分子的空间构型
                
                \cemh{HCN} \(2 + 0 = 2\)对，\textcolor{blue}{直线形}
            
                \cemh{HClO} \cemh{HOCl} \(2 + 2 = 4\)对，\textcolor{blue}{角形}
                
                \cemh{COCl2} \(3 + 0 = 3\)对，\textcolor{blue}{平面三角形}

                \textcolor{red}{斥力大小}：三键 \(>\) 双键 \(>\) 单键\quad \(\longleftarrow\)\textcolor{red}{为什么？}
            \end{xmp}

            知道分子的空间结构，能否帮助我们分析分子的成键方式？

            \begin{xmp}{空间构型练习}
                分析下列分子或离子的空间构型
                
                \cemh{SO3} \(3 + 0 = 3\)对，\textcolor{blue}{平面正三角形}
                
                \cemh{SO4^{2-}} \(4 + 0 = 4\)对，\textcolor{blue}{正四面体形}
                
                \cemh{NF3} \(3 + 1 = 4\)对，\textcolor{blue}{三角锥形}
            \end{xmp}

            \noindent\textbf{比较下列分子的键角大小}

            \makebox[0.5\textwidth][l]{\Circled{1} 对于\cemh{H2O}和\cemh{H2S}}\Circled{2} 对于\cemh{NH3}和\cemh{NF3}

            % transposed table inserted below
            \begin{center}
            \small
            \begin{tabular}{|c|c|c|c|c|}
                \hline
                \noindent\textbf{分子} & \cemh{H2O} & \cemh{H2S} & \cemh{NH3} & \cemh{NF3}\\
                \hline
                \noindent\textbf{键角} & \ang{104.5} & \ang{92.2} & \ang{107.3} & \ang{102}\\
                \hline
            \end{tabular}
            \end{center}

            一般来说：中心原子电负性越大，键角越大

            \phantom{一般来说：}配位原子电负性越大，键角越小

    \subsection{杂化轨道理论分析分子或复杂离子的结构}\label{subsec:杂化轨道理论}

        \noindent\textbf{归纳：}
        \begin{enumerate}
            \item 先由价层电子对数判断杂化方式。
            \item 几个配位原子形成几根 \(\upsigma\) 键。
            \item 结合VSEPR模型和配位原子数分析空间结构。
            \item sp$^2$杂化，分析垂直于成键原子平面的p轨道如何形成\(\uppi\)键。
            \item sp杂化，分析垂直于键轴的p轨道如何形成\(\uppi\)键。
            \item 若有多个原子的p轨道均垂直，则形成离域\(\uppi\)键。
        \end{enumerate}

        \textcolor{red}{一般不分析末端原子的杂化方式（大多数时候其实也杂化）。}

        \subsubsection{实例：直线形分子}\label{subsubsec:直线形分子}

            \noindent\begin{minipage}[t]{0.48\textwidth}
                \noindent\textbf{\cemh{BeCl2}}
                
                VSEPR模型：直线形
                
                Be原子 sp 杂化
                
                \(\therefore\) \cemh{BeCl2}直线形
                
                \vspace{\baselineskip}
            
                \noindent\textbf{\cemh{HCN}}
                
                VSEPR模型：直线形
                
                C原子 sp 杂化
                
                \cemh{H-C#N}
                
                \(\therefore\) \cemh{HCN}直线形
            \end{minipage}
            \hfill
            \begin{minipage}[t]{0.48\textwidth}
                \noindent\textbf{\cemh{CO2}}
                
                VSEPR模型：直线形
                
                C原子 sp 杂化
                
                \cemh{CO2}分子中含有2根\(\upsigma\)键、2根\(\uppi\)键
                
                \(\therefore\) \cemh{CO2}直线形
            \end{minipage}

        \subsubsection{实例：平面三角形分子}\label{subsubsec:平面三角形分子}

            \noindent\begin{minipage}[t]{0.48\textwidth}
                \noindent\textbf{\cemh{BF3}}
                
                VSEPR模型：平面正三角形
                
                B原子 sp$^2$ 杂化
                
                \(\therefore\) \cemh{BF3}平面正三角形
            \end{minipage}
            \hfill
            \begin{minipage}[t]{0.48\textwidth}
                \noindent\textbf{\cemh{CO3^{2-}}}
                
                VSEPR模型：平面正三角形
                
                C原子 sp$^2$ 杂化
            \end{minipage}

        \subsubsection{实例：四面体形分子}\label{subsubsec:四面体形分子}

            \noindent\begin{minipage}[t]{0.48\textwidth}
                \noindent\textbf{\cemh{CH4}、\cemh{CF4}}
                
                VSEPR模型：正四面体形
                
                C原子 sp$^3$ 杂化
            \end{minipage}
            \hfill
            \begin{minipage}[t]{0.48\textwidth}
                \noindent\textbf{\cemh{NH4+}}
                
                N原子 sp$^3$ 杂化
                
                \(\therefore\) \cemh{NH4+}正四面体形
            \end{minipage}

        \subsubsection{实例：三角锥形分子}\label{subsubsec:三角锥形分子}

            \noindent\textbf{\cemh{NH3}}

            VSEPR模型：四面体形

            N原子 sp$^3$ 杂化

            \(\therefore\) \cemh{NH3}三角锥形

        \subsubsection{实例：角形分子}\label{subsubsec:角形分子}

            \noindent\begin{minipage}[t]{0.48\textwidth}
                \noindent\textbf{\cemh{H2O}}
                
                VSEPR模型：四面体形
                
                O原子 sp$^3$ 杂化
                
                \(\therefore\) \cemh{H2O}角形
                
                \vspace{\baselineskip}
                
                \noindent\textbf{\cemh{H3O+}}
                
                \(\therefore\) \cemh{H3O+}三角锥形
            \end{minipage}
            \hfill
            \begin{minipage}[t]{0.48\textwidth}
                \noindent\textbf{\cemh{HOCl}}
                
                VSEPR模型：四面体形
                
                O原子 sp$^3$ 杂化
                
                \(\therefore\) \cemh{HOCl}角形
            \end{minipage}

        \subsubsection{实例：\texorpdfstring{\cemh{SO2}}{SO₂}和\texorpdfstring{\cemh{SO3}}{SO₃}}\label{subsubsec:SO2和SO3}

            \noindent\begin{minipage}[t]{0.48\textwidth}
                \noindent\textbf{\cemh{SO2}}
                
                VSEPR模型：平面三角形
                
                S原子 sp$^2$ 杂化
                
                \cemh{SO2}分子含有\(\uppi_3^4\)离域\(\uppi\)键
                
                \(\therefore\) \cemh{SO2}角形
            \end{minipage}
            \hfill
            \begin{minipage}[t]{0.48\textwidth}
                \noindent\textbf{\cemh{SO3}}
                
                VSEPR模型：平面正三角形
                
                S原子 sp$^2$ 杂化
                
                一个O原子空出一个2p轨道
                
                \cemh{SO3}分子含有\(\uppi_4^6\)离域\(\uppi\)键
                
                \(\therefore\) \cemh{SO3}平面正三角形
            \end{minipage}

    \subsection{等电子体原理}\label{subsec:等电子体原理}

        \noindent\textbf{等电子体：}原子数相同且价电子总数相同的分子或离子。

        具有相同的结构特征：化学键类型和空间结构。

        \begin{center}
        \small
        \begin{tabular}{|c|l|}
            \hline
            \noindent\textbf{等电子体类别} & \multicolumn{1}{c|}{\textbf{等电子体类别}} \\
            \hline
            3原子8电子 & \cemh{H2O}、\cemh{H2S}、\cemh{NH2-} \\
            \hline
            4原子8电子 & \cemh{NH3}、\cemh{PH3}、\cemh{H3O+}、\cemh{CH3-} \\
            \hline
            5原子8电子 & \cemh{CH4}、\cemh{NH4+} \\
            \hline
            2原子10电子 & \cemh{N2}、\cemh{CO}、\cemh{CN-}、\cemh{C2^{2-}}、\cemh{NO+} \\
            \hline
            3原子16电子 & \cemh{CO2}、\cemh{CS2}、\cemh{N2O}、\cemh{BeCl2}、\cemh{NO2+} \\
            \hline
            3原子18电子 & \cemh{SO2}、\cemh{O3}、\cemh{NO2-} \\
            \hline
            4原子24电子 & \cemh{BF3}、\cemh{CO3^{2-}}、\cemh{NO3-}、\cemh{SO3} \\
            \hline
            4原子26电子 & \cemh{PCl3}、\cemh{SO3^{2-}}、\cemh{ClO3^-}、\cemh{BrO3^-} \\
            \hline
            5原子32电子 & \cemh{CF4}、\cemh{CCl4}、\cemh{SO4^{2-}}、\cemh{ClO4-}、\cemh{PO4^{3-}} \\
            \hline
        \end{tabular}
        \end{center}

        \begin{itemize}
            \item \cemh{PF5}：sp$^3$d杂化，三角双锥
            \item \cemh{SF6}：sp$^3$d$^2$杂化，正八面体
        \end{itemize}
